\documentclass[11pt,a4paper]{article}
\usepackage[a4paper,margin=2.5cm]{geometry}
\setlength{\parindent}{0pt}
\setlength{\parskip}{0.6em}
\usepackage{helvet}
\renewcommand{\familydefault}{\sfdefault}
\usepackage[utf8]{inputenc}
\usepackage[T1]{fontenc}
\usepackage{setspace}
\onehalfspacing
\usepackage{titlesec}
\titleformat{\section}{\bfseries\Large\centering}{}{0pt}{}
\titlespacing*{\section}{0pt}{2em}{1em}

\begin{document}

\begin{center}
\textbf{\Large POLÍTICA FINANCEIRA}\\
\textbf{\large MOVIMENTO EM DEFESA DA EDUCAÇÃO -- MoveEduca}
\end{center}

\section*{1. FINALIDADE}

\textbf{1.1.} Esta Política Financeira estabelece regras simples para planejamento, recebimento, aplicação, autorização, pagamento, controle e prestação de contas dos recursos do MoveEduca.

\textbf{1.2.} A política complementa o Estatuto Social, o Regimento Interno e as demais políticas institucionais.

\section*{2. PRINCÍPIOS}

\textbf{2.1.} A gestão financeira observará legalidade, transparência, economicidade, segregação de funções, finalidade institucional, documentação adequada e prestação de contas.

\textbf{2.2.} Todo recurso será aplicado integralmente nas finalidades estatutárias, vedada distribuição de lucros, sobras, excedentes, vantagens indevidas ou benefícios pessoais injustificados.

\section*{3. ORÇAMENTO E PLANEJAMENTO}

\textbf{3.1.} A Diretoria Executiva elaborará orçamento ou previsão financeira anual, ainda que simplificada, contendo estimativa de receitas, despesas, projetos e necessidades operacionais.

\textbf{3.2.} Projetos com orçamento próprio deverão ter registro mínimo de fonte de recursos, finalidade, responsável, período de execução, despesas previstas e forma de prestação de contas.

\section*{4. RECEITAS}

\textbf{4.1.} O MoveEduca poderá receber contribuições, doações, subvenções, patrocínios, apoios, recursos de editais, termos de colaboração, termos de fomento, acordos de cooperação, contratos, prestação de serviços, produtos institucionais, rendimentos financeiros e outras receitas lícitas compatíveis com o Estatuto Social.

\textbf{4.2.} Doações relevantes ou com finalidade específica deverão ser documentadas por termo, recibo, comprovante bancário ou instrumento equivalente.

\textbf{4.3.} A associação poderá recusar recursos que imponham condição ilícita, vantagem indevida, conflito com sua missão, risco reputacional relevante ou desvio de finalidade.

\section*{5. DESPESAS E ALÇADAS}

\textbf{5.1.} Despesas deverão ter finalidade institucional, previsão orçamentária ou justificativa, comprovação documental e autorização compatível com o valor e a natureza do gasto.

\textbf{5.2.} Despesas ordinárias, recorrentes ou de baixo valor poderão ser autorizadas pelo Presidente ou pelo Diretor Financeiro, desde que registradas e comprovadas.

\textbf{5.3.} Despesas relevantes, extraordinárias ou não previstas deverão ser aprovadas pela Diretoria Executiva.

\textbf{5.4.} Operações patrimoniais relevantes, aquisição, alienação ou oneração de imóveis observarão o Estatuto Social e dependerão das aprovações nele previstas.

\textbf{5.5.} É vedado fracionar despesas para evitar aprovações, cotações, controles ou prestação de contas.

\section*{6. PAGAMENTOS E MOVIMENTAÇÃO BANCÁRIA}

\textbf{6.1.} A movimentação bancária e a autorização de pagamentos observarão assinatura ou autorização conjunta do Presidente e do Diretor Financeiro, conforme o Estatuto Social.

\textbf{6.2.} Pagamentos deverão ser realizados preferencialmente por meio bancário rastreável.

\textbf{6.3.} Saques em espécie deverão ser excepcionais, justificados, comprovados e registrados.

\textbf{6.4.} Adiantamentos deverão ter finalidade definida, responsável identificado, prazo de comprovação e devolução de saldo não utilizado.

\section*{7. COMPRAS, CONTRATAÇÕES E REEMBOLSOS}

\textbf{7.1.} Compras e contratações deverão buscar compatibilidade com valores de mercado, qualidade adequada, documentação e ausência de conflito de interesses.

\textbf{7.2.} Sempre que razoável, deverão ser obtidas referências de preço ou cotações, especialmente em despesas relevantes.

\textbf{7.3.} Reembolsos somente serão pagos mediante autorização, vínculo com atividade institucional e apresentação de comprovantes.

\textbf{7.4.} Despesas sem comprovante idôneo poderão ser recusadas, salvo hipóteses excepcionais justificadas pela Diretoria Executiva.

\section*{8. RECURSOS PÚBLICOS E PROJETOS VINCULADOS}

\textbf{8.1.} Recursos públicos serão aplicados conforme a legislação, o edital, o plano de trabalho, o instrumento firmado e as orientações do órgão parceiro.

\textbf{8.2.} Quando exigido, recursos de projetos serão controlados separadamente por conta bancária, centro de custo, planilha, relatório ou outro meio adequado.

\textbf{8.3.} Despesas incompatíveis com plano de trabalho, rubrica aprovada ou instrumento de parceria não poderão ser realizadas com recursos públicos.

\section*{9. DOCUMENTAÇÃO, CONTABILIDADE E PRESTAÇÃO DE CONTAS}

\textbf{9.1.} A Diretoria Executiva manterá documentos financeiros, fiscais, bancários, contábeis e contratuais organizados e disponíveis para análise dos órgãos competentes.

\textbf{9.2.} A escrituração contábil observará as Normas Brasileiras de Contabilidade aplicáveis às entidades sem fins lucrativos.

\textbf{9.3.} As demonstrações financeiras e prestações de contas serão submetidas ao Conselho Fiscal e à Assembleia Geral, conforme o Estatuto Social.

\textbf{9.4.} Documentos financeiros deverão ser guardados pelo prazo legal aplicável ou pelo prazo maior exigido por instrumento de parceria, edital, financiador ou órgão público.

\section*{10. CONFLITO DE INTERESSES}

\textbf{10.1.} Qualquer pessoa envolvida em decisão financeira deverá declarar conflito de interesses real, potencial ou aparente.

\textbf{10.2.} A pessoa em conflito não poderá aprovar, autorizar, fiscalizar isoladamente ou influenciar a despesa ou contratação correspondente.

\section*{11. DISPOSIÇÕES FINAIS}

\textbf{11.1.} Casos omissos serão decididos pela Diretoria Executiva, com comunicação ao Conselho Fiscal quando envolver risco financeiro, contábil ou patrimonial relevante.

\textbf{11.2.} Esta Política Financeira poderá ser revisada sempre que necessário para adequação à legislação, ao Estatuto Social, às parcerias e à realidade operacional do MoveEduca.

\vspace{2cm}
\begin{center}
\textbf{[CIDADE/UF]}, \textbf{[DATA]}.

\vspace{2cm}
\rule{8cm}{0.4pt}\\
\textbf{[NOME COMPLETO]}\\
Presidente do MoveEduca

\vspace{1.5cm}
\rule{8cm}{0.4pt}\\
\textbf{[NOME COMPLETO]}\\
Diretor(a) Financeiro(a) do MoveEduca

\vspace{1.5cm}
\rule{8cm}{0.4pt}\\
\textbf{[NOME DO ADVOGADO]}\\
\textbf{[OAB/UF Nº]}
\end{center}

\end{document}
